\chapter[RESULTADOS]{RESULTADOS}

Este capítulo reporta os resultados obtidos até o momento. Ele começa pelas verificações que sustentam o método descrito no Capítulo 3 --- e que são pré-requisito para que os experimentos façam sentido --- e prossegue para os resultados dos experimentos propriamente ditos.

% =====================================================================
% TODO (estrutura): este capítulo está completo apenas nas seções 4.1 a 4.4.
% As seções do Experimento 1 e do Experimento 2 aguardam execução. Ver
% docs/roadmap.md no repositório para o estado corrente.
% =====================================================================

\section{Verificação da Equivalência da Varredura Fora de Linha}

A Seção \ref{sec:varredura} afirma que executar o Professor uma única vez com o limiar no piso e filtrar o material armazenado é exatamente equivalente a reexecutar o modelo em cada limiar. Como todo o protocolo de calibração depende dessa afirmação, ela foi verificada por execução, e não apenas por inspeção do código.

O procedimento fixou uma imagem e, para cada uma de quatro sondas, comparou duas grandezas em oito limiares distribuídos entre 0,005 e 0,500: o resultado de executar o Professor com aquele limiar e o resultado de filtrar, naquele mesmo valor, o material capturado com o limiar no piso de 0,001. A comparação abrangeu o número de detecções, os escores e os vértices dos polígonos, \textbf{sem tolerância numérica}.

As 32 comparações resultaram idênticas em todas as três grandezas. Não houve divergência em nenhum limiar, nem no número de regiões, nem em qualquer coordenada de polígono. A Tabela \ref{tab:d18} resume o número de detecções sobreviventes em cada configuração --- valor que, por ser idêntico nos dois métodos, é reportado em coluna única.

\begin{table}[ht]
  \caption{Detecções sobreviventes por limiar, idênticas entre execução real e varredura fora de linha.}
  \label{tab:d18}
  \centering
  \begin{tabular}{lrrrrrrrr}
    \hline
    \bfseries Sonda & \bfseries 0,005 & \bfseries 0,02 & \bfseries 0,05 & \bfseries 0,08 & \bfseries 0,12 & \bfseries 0,20 & \bfseries 0,35 & \bfseries 0,50 \\
    \hline
    \textit{crack}        & 168 & 165 & 118 & 76 & 41 & 20 & 9 & 4 \\
    \textit{vein}         & 164 & 149 & 71  & 34 & 21 & 11 & 7 & 4 \\
    \textit{Stain}        & 195 & 186 & 132 & 94 & 74 & 48 & 15 & 7 \\
    \textit{Dark patches} & 194 & 175 & 120 & 73 & 52 & 25 & 9 & 6 \\
    \hline
  \end{tabular}
  \legend{Fonte: elaborado pelo autor.}
\end{table}

O resultado tem duas consequências. A primeira é prática: o ciclo de calibração deixa de custar minutos de GPU por iteração e passa a custar milissegundos, o que é o que viabiliza submeter um mesmo limiar às quatro imagens da litologia simultaneamente. A segunda é metodológica: a varredura fora de linha não introduz aproximação alguma, de modo que nenhum resultado obtido por ela precisa ser qualificado como estimativa.

Ressalva necessária: a equivalência é propriedade da \textbf{implementação} verificada, não do método em abstrato. Uma versão futura da biblioteca que aplicasse o limiar antes da supressão de não-máximos, ou que o incorporasse ao próprio modelo, invalidaria a conclusão. Por isso o procedimento de verificação integra o repositório e deve ser reexecutado a cada atualização de versão.

\section{Caracterização dos Escores do Professor}

A escolha da escala logarítmica na regra de limiar (Seção \ref{sec:regra}) decorre da distribuição observada dos escores, medida sobre a imagem de descoberta de uma litologia da faixa A e reportada na Tabela \ref{tab:escores}.

\begin{table}[ht]
  \caption{Distribuição dos escores do Professor por sonda, com o limiar no piso.}
  \label{tab:escores}
  \centering
  \begin{tabular}{lrrrr}
    \hline
    \bfseries Sonda & \bfseries Detecções & \bfseries Mínimo & \bfseries Mediana & \bfseries Máximo \\
    \hline
    \textit{crack}        & 168 & 0,016 & 0,074 & 0,699 \\
    \textit{vein}         & 164 & 0,013 & 0,045 & 0,694 \\
    \textit{Stain}        & 195 & 0,012 & 0,076 & 0,683 \\
    \textit{Dark patches} & 194 & 0,011 & 0,061 & 0,766 \\
    \hline
  \end{tabular}
  \legend{Fonte: elaborado pelo autor.}
\end{table}

Os escores percorrem quase duas ordens de grandeza, de 0,011 a 0,766, mas a mediana situa-se entre 0,045 e 0,076 em todas as sondas: a massa da distribuição concentra-se na década baixa e a cauda superior é rala. Em uma grade linear de limiares, portanto, a quase totalidade dos pontos cairia em região sem detecção alguma, e o joelho da curva recairia sistematicamente sobre seu primeiro ponto --- daí a normalização logarítmica do eixo.

O dado também recomenda cautela na transposição de valores da literatura de pseudo-rotulagem. Os limiares elevados ali reportados referem-se a escores de detectores supervisionados calibrados, cuja escala não corresponde à observada aqui. O que se transporta daquela literatura é a direção do compromisso --- privilegiar precisão --- e não a magnitude do corte.

\section{Geometria dos Polígonos Produzidos pelo Professor}

A conversão das máscaras do Professor em polígonos no formato de anotação apresenta uma particularidade com consequência sobre a qualidade do rótulo. Quando a máscara de uma detecção possui mais de um contorno desconexo, a implementação utilizada funde todos eles em uma única poligonal, ligando-os por trajetos de ida e volta.

A medição sobre a imagem de descoberta de uma litologia da faixa A, com a sonda \textit{crack} no limiar de 0,08, encontrou 45 detecções com mais de um contorno em um total de 76. Como os trajetos de ligação têm área praticamente nula, o efeito agregado é pequeno: a área total dos polígonos supera a das máscaras em 2,4\%, e a sobreposição entre a união das máscaras e a união dos polígonos atinge IoU de 0,928. O caso individual, porém, não é igualmente benigno: em uma detecção composta por 18 contornos, a área do polígono alcançou 2,5 vezes a da máscara correspondente.

A implicação é que o efeito é desprezível quando se avalia a cobertura agregada, mas não quando uma detecção fragmentada é tomada isoladamente como instância de treinamento. Os trajetos de ligação também se mostraram visíveis na prática: durante a calibração, cada um aparecia como uma reta atravessando a chapa.

Adotou-se, por isso, a separação: \textbf{cada contorno da máscara torna-se um polígono próprio}, sem trajeto de ligação, tanto no material de calibração quanto nos arquivos de anotação. A alternativa de reter apenas o maior contorno foi descartada porque elimina área efetivamente marcada pelo Professor. A separação aumenta o número de polígonos por imagem, mas não altera a unidade sobre a qual a regra de limiar opera: o escore pertence à detecção, e os contornos de uma mesma detecção entram e saem juntos na varredura. Uma marcação, na curva da Seção \ref{sec:regra}, continua sendo uma detecção do Professor. Verificou-se, nas litologias calibradas à época, que as contagens por imagem no limiar adotado e o limiar produzido pela regra permaneceram idênticos após a mudança.

Duas questões permanecem em aberto no pós-processamento: a área mínima abaixo da qual um fragmento é descartado, e se a curva da regra deveria contar contornos em vez de detecções. Esta segunda questão só surgiu depois de a forma da regra ter sido fixada; caso seja adotada, os resultados serão reportados nas duas versões, com a indicação de que a segunda foi definida posteriormente.

\section{Classificador de Litologias}
\label{sec:res_classificador}

A reprodução do classificador (Seção \ref{sec:classificador}) atingiu \textbf{96,72\%} de acurácia sobre as 5.153 imagens da partição de teste original, com acurácia balanceada de 92,32\% e F1 macro de 0,924. Os erros concentram-se em poucas litologias --- 15 das 45 apresentam algum --- e em poucos pares: \textit{naica} confundida com \textit{white\_olympus} em 33 imagens, \textit{white\_himalaya} com \textit{white\_liberdade} em 31, e \textit{white\_superiore} com \textit{kalahari} em 21 e com \textit{white\_olympus} em 15.

O caso da \textit{white\_superiore} indica a causa. O modelo acertou 35 de suas 37 imagens de validação e nenhuma de suas 36 imagens de teste. Na inspeção visual, as imagens de teste são acinzentadas, com veios lineares e paralelos, enquanto as de treinamento e de validação são brancas, com veios ramificados: as de teste aparentam pertencer a outro bloco. Uma sondagem por similaridade de cor e brilho reforça essa leitura para essa litologia, mas não se confirma em todas --- em outras, o teste está tão próximo do treinamento quanto a validação.

Para isolar o efeito do particionamento, o mesmo modelo foi treinado com a mesma receita e as partições sorteadas por imagem, como no trabalho original. O resultado foi de \textbf{99,92\%} de acurácia, com 4 erros em 5.195 imagens, acima dos 99,21\% reportados pelo DeepStoneAI. A diferença de cerca de três pontos vem, portanto, da divisão dos dados, e não da resolução nem do modelo: no sorteio por imagem, chapas vizinhas das de treinamento --- cortadas do mesmo bloco e com os mesmos veios --- caem na partição de teste.

Os dois números descrevem cenários de uso distintos, e nenhum é tratado como erro do outro. O de 99,92\% corresponde a classificar chapas de um bloco já visto, situação comum na indústria, em que um mesmo bloco gera dezenas de chapas. O de 96,72\% aproxima-se de classificar chapas de um bloco novo. O valor do DeepStoneAI não é diretamente comparável a nenhum dos dois, por adotar particionamento próprio. O modelo empregado no \textit{pipeline} é o treinado sobre as partições originais.

\section{Experimento 1 --- Critério por Litologia \textit{versus} Critério Único}

% TODO: executar e redigir. Depende da calibração das litologias da faixa A e da
% anotação do conjunto de referência (Seção \ref{sec:referencia}).
% Reportar, pelos quatro braços e por litologia: IoU, precisão, recall e taxa de
% falso positivo; incluir o par de figuras "critério único x critério calibrado"
% sobre a mesma chapa; e reportar a taxa de concordância entre a regra e o autor.
% Figuras de exemplo (decisão do Henrique, 10/10/2026): o lado calibrado usa o
% LIMIAR DE TRABALHO (braço "calibrado pelo autor"), a não ser que o joelho da
% regra produza figura melhor. A legenda diz qual dos dois foi usado.

\section{Experimento 2 --- Especialistas \textit{versus} Generalista}

% TODO: executar e redigir, faixa por faixa. Depende da geração do conjunto de
% treinamento pelo Professor, do pós-processamento discutido na Seção 4.3 e do
% treinamento dos três braços descritos no Quadro \ref{quad:bracos_exp2}.
